% ═══════════════════════════════════════════════════════════════
%  PART XV — APRIL 8, 2026: THE MASTER IDENTITY AND GRAPH RH
% ═══════════════════════════════════════════════════════════════
\clearpage
\part{The Master Identity: Resolvent, Ihara Zeta, and the Graph Riemann Hypothesis}

\section{The Resolvent Identity}

The resolvent $R(x) = \Tr(xI-A)^{-1}$ of the W(3,3) adjacency
matrix is $R(x) = \frac{1}{x-k} + \frac{f}{x-r} + \frac{g}{x-s}$.

\begin{theorem}[The Spectral-Cyclotomic Bridge]
\label{thm:resolvent}
At $x=-1$, the eigenvalue denominators become cyclotomic values:
\begin{equation}
-1-k = -\Phithree, \qquad -1-r = -q, \qquad -1+|s| = q.
\end{equation}
Therefore:
\begin{equation}
\boxed{R(-1) = -\frac{1}{\Phithree} + \frac{g-f}{q} = -\frac{1}{\Phithree} - q = -\frac{v}{\Phithree},}
\end{equation}
which is equivalent to $1 + q\Phithree = v$---the vertex count formula
$v = (q^4{-}1)/(q{-}1)$. The shift $x\mapsto -(x{+}1)$ bridges
spectral graph theory and cyclotomic number theory.
\end{theorem}


\section{The Ihara Zeta Function and Graph Riemann Hypothesis}

The Ihara zeta function of W(3,3) has non-trivial factors:
\begin{align}
\text{$r$-factor:}& \quad 1-2u+11u^2, \quad \Delta_r = 4-44 = -v = -40,\\
\text{$s$-factor:}& \quad 1+4u+11u^2, \quad \Delta_s = 16-44 = -(v{-}k) = -28.
\end{align}

\begin{theorem}[Graph Riemann Hypothesis for W(3,3)]
\label{thm:graph-rh}
All non-trivial Ihara zeta zeros of W(3,3) satisfy
\begin{equation}
\boxed{|u| = \frac{1}{\sqrt{k-1}} = \frac{1}{\sqrt{11}}.}
\end{equation}
Both eigenvalues satisfy $|r|, |s| \leq 2\sqrt{k-1}$,
so W(3,3) is a \textbf{Ramanujan graph}.
Its Ihara zeta function satisfies the graph-theoretic Riemann Hypothesis:
all non-trivial zeros lie on the ``critical circle'' $|u|=(k{-}1)^{-1/2}$,
analogous to the classical RH assertion $\Re(s)=\tfrac{1}{2}$.
\end{theorem}

\begin{corollary}[Ihara Discriminants]
The discriminant of the $r$-factor is $\Delta_r = -v$;
the discriminant of the $s$-factor is $\Delta_s = -(v{-}k)$.
The vertex count and the complement degree appear directly
as Ihara discriminants.
\end{corollary}

\section{Trace Formulas}

The moments $\Tr(A^n) = k^n + f\cdot r^n + g\cdot s^n$ encode:
\begin{center}
\begin{tabular}{@{}clll@{}}
\toprule
$n$ & $\Tr(A^n)$ & Identity & Meaning \\
\midrule
0 & 40 & $v$ & vertex count \\
1 & 0 & $k+fr+gs$ & traceless (critical for physics) \\
2 & 480 & $vk = 2E = a_0$ & spectral action coefficient \\
3 & 960 & $6\mu v$ & number of triangles $= \mu v = 160$ \\
\bottomrule
\end{tabular}
\end{center}
The tracelessness $\Tr(A)=0$ is equivalent to $g-f = -q^2$, which
forces the multiplicities once the eigenvalues are known.
\end{corollary}
